\section{Ordinary convergent reconstruction and polygamma identities}
\label{sec:ordinary-reconstruction}

The rank theorem is more than a statement about a determinant.  It
allows every late zeta coordinate to be extracted by a finite rational
combination of centered moments.  The analytic continuation can then be
converted into an ordinary absolutely convergent series by subtracting
only the nonnegative powers in the centered asymptotic expansion.

\subsection{Centered asymptotic coefficients}

For an integer $s\ge2$ and $x=n+1/2$, the shifted Euler--Maclaurin
expansion~\cite{DLMF} is
\begin{equation}
 A_s(n)\sim\sum_{m\ge0}\alpha_{s,m}x^{1-s-2m},
 \label{eq:centered-tail-asymptotic}
\end{equation}
where
\begin{equation}
 \alpha_{s,0}=\frac1{s-1},
 \qquad
 \alpha_{s,m}=\frac{B_{2m}(1/2)}{(2m)!}(s)_{2m-1}
 \quad(m\ge1).
 \label{eq:alpha}
\end{equation}
Here $(s)_j=s(s+1)\cdots(s+j-1)$ and
$B_{2m}(1/2)=(2^{1-2m}-1)B_{2m}$.  The half-integer center removes all
odd corrections.  For the order-ten pair define
\begin{equation}
 \beta_{b,\ell}=\sum_{u=0}^{\ell}
       \alpha_{10,u}\alpha_{2b,\ell-u}.
 \label{eq:beta}
\end{equation}
Then
\begin{equation}
 A_{10}(n)A_{2b}(n)
 \sim\sum_{\ell\ge0}\beta_{b,\ell}x_n^{-2b-8-2\ell}.
 \label{eq:product-asymptotic}
\end{equation}

\subsection{The rational extraction rows}

Let $N^{(b)}=M_{5,b}^{-1}$, which exists by
\cref{thm:complete-order-ten}.  For $0\le k<b$ and
$0\le r\le2b-1$, set
\begin{equation}
 w_{b,k,r}=\sum_{t=0}^{b-1}N^{(b)}_{kt}
      \binom{2(b+t)+1}{2r}
      2^{-2(b+t-r)},
 \label{eq:reconstruction-weight}
\end{equation}
where the binomial coefficient is understood to be zero when
$r>b+t$.  Define the even polynomial
\begin{equation}
 P_{b,k}(x)=\sum_{r=0}^{2b-1}w_{b,k,r}x^{2r}.
 \label{eq:Pbk}
\end{equation}
Its nonnegative asymptotic part against the product of tails is
\begin{equation}
 Q_{b,k}(x)=
 \sum_{r=0}^{2b-1}w_{b,k,r}
 \sum_{\substack{\ell\ge0\\2r-2b-8-2\ell\ge0}}
     \beta_{b,\ell}x^{2r-2b-8-2\ell}.
 \label{eq:Qbk}
\end{equation}
Every sum is finite and rational.  In particular,
\begin{equation}
 \deg P_{b,k}\le4b-2,
 \qquad
 \deg Q_{b,k}\le2b-10.
 \label{eq:PQ-degrees}
\end{equation}

\begin{theorem}[All-coordinate ordinary reconstruction]
\label{thm:ordinary-reconstruction}
For every $b\ge6$ and $0\le k<b$,
\begin{equation}
 \boxed{\displaystyle
 \sum_{n=0}^{\infty}
 \left\{
 P_{b,k}(x_n)A_{10}(n)A_{2b}(n)-Q_{b,k}(x_n)
 \right\}=\zeta(2k),}
 \qquad x_n=n+\frac12,
 \label{eq:ordinary-reconstruction}
\end{equation}
where $\zeta(0)=-1/2$.  The series is absolutely convergent, with
summand $O(n^{-2})$.
\end{theorem}

\begin{proof}
From \eqref{eq:late-coordinates},
\[
 Z_k=\sum_{t=0}^{b-1}N^{(b)}_{kt}C_{b+t}.
\]
Substitution of \eqref{eq:C-transform} gives
\[
 Z_k=\sum_{r=0}^{2b-1}w_{b,k,r}D_r.
\]
By definition, the right side is the centered regularized sum of
$P_{b,k}(x_n)A_{10}(n)A_{2b}(n)$.  Equations
\eqref{eq:product-asymptotic} and \eqref{eq:Qbk} show that subtraction
of $Q_{b,k}$ removes every nonnegative power, leaving $O(x_n^{-2})$.
Hence the subtracted series converges absolutely.

It remains to compare the ordinary sum with the centered regularized
one.  Every monomial in $Q_{b,k}$ has nonnegative even degree, and
\[
 \sum_{n\ge0}^{\mathrm{a.c.}}(n+\tfrac12)^{2j}
   =\zeta(-2j,\tfrac12)=0\qquad(j\ge0).
\]
Thus the finite subtraction has zero centered zeta-regularized sum.
The ordinary subtracted series therefore equals the inherited regular
value $\sum_rw_{b,k,r}D_r=Z_k=\zeta(2k)$.
\end{proof}

\begin{corollary}[Product-polygamma form]
\label{cor:polygamma-reconstruction}
For every $b\ge6$ and $0\le k<b$,
\begin{equation}
 \sum_{n=0}^{\infty}
 \left\{
 \frac{P_{b,k}(x_n)}{9!(2b-1)!}
 \psi^{(9)}(n+1)\psi^{(2b-1)}(n+1)
 -Q_{b,k}(x_n)
 \right\}=\zeta(2k).
 \label{eq:polygamma-reconstruction}
\end{equation}
\end{corollary}

\begin{proof}
For even $s$, equation \eqref{eq:polygamma-tail} reads
$A_s(n)=\psi^{(s-1)}(n+1)/(s-1)!$.  Substitute this into
\eqref{eq:ordinary-reconstruction}.
\end{proof}

\begin{remark}[Effective, not merely existential]
All coefficients in \eqref{eq:ordinary-reconstruction} are obtained by
finite rational arithmetic: form \eqref{eq:source-M}, invert it, apply
\eqref{eq:reconstruction-weight}, and truncate the explicit asymptotic
convolution \eqref{eq:Qbk}.  The accompanying identity generator follows
exactly this procedure.  No recognition of numerical constants or fit
to a floating-point series is involved.
\end{remark}
